\chapter{Đổi Chỗ Để Tám}
\chapterintrobi{Lý do đổi chỗ rất học đường, mục đích thì khá xã hội}{Change seats to chat}{Đôi khi chiến lược ngồi gần nhau còn được chuẩn bị kỹ hơn chiến lược làm bài.}

\begin{lucbatbox}{Lục bát: Ngồi gần cho tiện chuyện}
Em xin đổi chỗ nhẹ nhàng\
Nói là dễ học, dễ bàn, dễ trao\
Nhưng vào tiết mới biết sao\
Trao là trao chuyện, còn bài để sau\
Bạn cười, em cũng cúi đầu\
Cả hai hợp tác làm đau sự im
\end{lucbatbox}

\begin{tutuyetbox}{Thất ngôn tứ tuyệt: Đồng đội đúng chỗ}
Đổi chỗ nghe như chuyện học hành\
Vào ngồi mới rõ ý em nhanh\
Nếu phần bàn bài chăm như chuyện nhỏ\
Hai người điểm số hẳn long lanh
\end{tutuyetbox}

\begin{ghichu}
Đổi chỗ để tám là kiểu đầu tư rất hiệu quả cho giải trí và khá rủi ro cho sự tập trung.
\end{ghichu}

\begin{englishnote}
\textbf{change seats}: đổi chỗ.\par
\textbf{chat in class}: nói chuyện trong lớp.\par
\textbf{We just want to sit together.}: Bọn em chỉ muốn ngồi gần nhau thôi.
\end{englishnote}